\documentclass{article}
\usepackage[utf8]{inputenc}
\usepackage[french]{babel}
\usepackage{booktabs}
\usepackage{xcolor}
\usepackage{amsmath}
\usepackage{siunitx}
\usepackage{chemfig}
\usepackage{chemformula}
\usepackage[version=4]{mhchem}

\DeclareSIUnit{\nothing}{\relax}
\graphicspath{ {./images/} }

\title{Schéma de Lewis et Théorie VSEPR}
\author{Yannick Patois}
\date{14/01/2026}

\begin{document}

\maketitle

\section{Schéma de Lewis}
\subsection{Introduction}
Les schémas de Lewis représentent les électrons de valence des atomes et des molécules.
Ils permettent de visualiser les liaisons covalentes, les paires libres et les charges formelles,
ce qui est essentiel pour comprendre la réactivité chimique et la géométrie moléculaire.

\subsection{Règles de Construction}
Pour dessiner un schéma de Lewis~:
\begin{enumerate}
    \item Déterminer le nombre total d'électrons de valence de tous les atomes dans la molécule.
    \item Ajouter ou retirer des électrons si la molécule est un ion.
    \item Placer les atomes en respectant leur connectivité (l'atome le moins électronégatif
    est généralement central).
    \item Distribuer les électrons pour former des liaisons simples entre les atomes.
    \item Compléter les octets (ou duets pour l'hydrogène) en ajoutant des paires libres ou
    des liaisons multiples si nécessaire.
    \item Vérifier les charges formelles et les indiquer si besoin.
\end{enumerate}

\subsection{Exemples de Schémas de Lewis}
\subsubsection{Molécules Simples}
\begin{itemize}
    \item \textbf{Eau (\ch{H2O})}~:
    L'oxygène (6 électrons de valence) forme deux liaisons simples avec deux hydrogènes
    (1 électron chacun). Il reste deux paires libres sur l'oxygène.
    \begin{center}
    \chemfig{H-[:52.24]\charge{45=\:,135=\:}{O}-[::-104.48]H}
    \end{center}

    \item \textbf{Dioxyde de carbone (\ch{CO2})}~:
    Le carbone (4 électrons de valence) forme deux doubles liaisons avec deux oxygènes
    (6 électrons chacun). Aucune paire libre sur le carbone.
    \begin{center}
        \chemfig{
        \charge{135=\:,225=\:}{O}
        =C
        =\charge{315=\:,45=\:}{O}
        }
    \end{center}

    \item \textbf{Ammoniac (\ch{NH3})}~:
    L'azote (5 électrons de valence) forme trois liaisons simples avec trois hydrogènes, laissant une paire libre.
    \begin{center}
        \chemfig{N(-[:180]\charge{135=\:,225=\:}{H})(-[:90]\charge{270=\:,330=\:}{H3})-[:0]\charge{30=\:,95=\:}{H4}}
    \end{center}
\end{itemize}

\subsubsection{Molécules avec Charges Formelles}
\begin{itemize}
    \item \textbf{Ion ammonium (\ch{NH4+})}~:
    L'azote forme quatre liaisons simples avec quatre hydrogènes, ce qui lui donne une charge formelle de +1.
    \begin{center}
    \chemfig{@N([:90]\charge{90=\:,150=\:}{H})(-[:150]\charge{30=\:,90=\:}{H})(-[:270]\charge{270=\:,330=\:}{H})-[:30]\charge{30=\:,95=\:}{H}}
    \end{center}

    \item \textbf{Ion carbonate (\ch{CO3^2-})}~:
    Le carbone forme trois doubles liaisons avec trois oxygènes, et porte une charge formelle de -2 répartie sur les oxygènes.
    \begin{center}
    \chemfig{@O-[:90]C(=[:150]\charge{150=\:,210=\:}{O^-})(=[:270]\charge{270=\:,330=\:}{O^-})=[:30]\charge{30=\:,90=\:}{O^-}}
    \end{center}
\end{itemize}

\subsection{Exceptions aux Règles de l'Octet}
Certaines molécules ne respectent pas la règle de l'octet~:
\begin{itemize}
    \item \textbf{Molécules hypovalentes}~: Comme \ch{BeCl2} ou \ch{BF3}, où l'atome central a moins de 8 électrons.
    \item \textbf{Molécules hypervalentes}~: Comme \ch{PCl5} ou \ch{SF6}, où l'atome central a plus de 8 électrons.
    \item \textbf{Radicaux libres}~: Comme \ch{NO} ou \ch{ClO2}, où un électron est non apparié.
\end{itemize}

\subsection{Application à la Théorie VSEPR}
Les schémas de Lewis sont une étape préliminaire à l'application de la théorie VSEPR. Ils permettent de~:
\begin{itemize}
    \item Identifier le nombre de paires liantes et non liantes autour de l'atome central.
    \item Déterminer la géométrie électronique et moléculaire.
    \item Prédire les angles de liaison et la polarité des molécules.
\end{itemize}

\subsection{Exercices d'Application}
\begin{enumerate}
    \item Dessinez le schéma de Lewis pour les molécules suivantes~: \ch{SO2}, \ch{SO3}, \ch{PO4^3-}.
    \item Identifiez les charges formelles dans \ch{HNO3} et \ch{ClO4-}.
    \item Proposez une géométrie moléculaire pour \ch{SF4} et \ch{XeF2} en utilisant la théorie VSEPR.
\end{enumerate}


\section{Théorie VSEPR et Géométrie Moléculaire}
\subsection{Principe de la Théorie VSEPR}
La théorie VSEPR (Valence Shell Electron Pair Repulsion) permet de prédire la 
géométrie des molécules en minimisant les répulsions entre les paires d'électrons de valence.

\subsection{Géométries Courantes}
\begin{center}
\begin{tabular}{>{\bfseries}l l l}
\toprule
Géométrie & Description & Exemple \\
\midrule
Linéaire & Angle de $180^\circ$ & \ce{CO2} \\
Trigonale plane & Angle de $120^\circ$ & \ce{BF3} \\
Tétraédrique & Angle de $109.5^\circ$ & \ce{CH4} \\
\bottomrule
\end{tabular}
\end{center}


\end{document}

